\documentclass[11pt]{article}
\newcommand{\ListaTitulo}{Gabarito --- Lista 02}
% Preâmbulo compartilhado — MATD48, Listas de Exercícios 2026
% Usado por Lista{NN}.tex e Gabarito{NN}.tex via % Preâmbulo compartilhado — MATD48, Listas de Exercícios 2026
% Usado por Lista{NN}.tex e Gabarito{NN}.tex via % Preâmbulo compartilhado — MATD48, Listas de Exercícios 2026
% Usado por Lista{NN}.tex e Gabarito{NN}.tex via \input{preamble}

\usepackage[utf8]{inputenc}
\usepackage[T1]{fontenc}
\usepackage[brazil]{babel}
\usepackage{fullpage}
\usepackage{amsmath,amssymb}
\usepackage{enumitem}
\usepackage{xcolor}
\usepackage{fancyhdr}
\usepackage{titlesec}
\usepackage{listings}
\usepackage{hyperref}
\usepackage{booktabs}
\usepackage{graphicx}

\definecolor{matd48azul}{HTML}{0B4F6C}
\definecolor{matd48verde}{HTML}{2E8B57}

\titleformat{\section}{\normalfont\Large\bfseries\color{matd48azul}}{\thesection}{1em}{}
\titleformat{\subsection}{\normalfont\large\bfseries\color{matd48azul}}{\thesubsection}{1em}{}

\providecommand{\ListaTitulo}{Lista de Exercícios}

\setlength{\headheight}{14pt}
\pagestyle{fancy}
\fancyhf{}
\lhead{\small MATD48 --- Planejamento de Experimentos A}
\rhead{\small \ListaTitulo}
\cfoot{\thepage}
\renewcommand{\headrulewidth}{0.4pt}

\lstset{
  basicstyle=\ttfamily\small,
  breaklines=true,
  frame=single,
  columns=fullflexible,
  backgroundcolor=\color{gray!8},
  commentstyle=\color{matd48verde},
  keywordstyle=\color{matd48azul}\bfseries,
  language=R,
  extendedchars=true,
  literate=%
    {á}{{\'a}}1 {à}{{\`a}}1 {â}{{\^a}}1 {ã}{{\~a}}1
    {é}{{\'e}}1 {ê}{{\^e}}1 {í}{{\'i}}1 {ó}{{\'o}}1
    {ô}{{\^o}}1 {õ}{{\~o}}1 {ú}{{\'u}}1 {ç}{{\c c}}1
    {Á}{{\'A}}1 {À}{{\`A}}1 {Â}{{\^A}}1 {Ã}{{\~A}}1
    {É}{{\'E}}1 {Ê}{{\^E}}1 {Í}{{\'I}}1 {Ó}{{\'O}}1
    {Ô}{{\^O}}1 {Õ}{{\~O}}1 {Ú}{{\'U}}1 {Ç}{{\c C}}1
}

\hypersetup{colorlinks=true, linkcolor=matd48azul, urlcolor=matd48azul, citecolor=matd48azul}

\newcommand{\cabecalho}[2]{%
  \begin{center}
    {\Large\bfseries\color{matd48azul} #1}\\[0.3em]
    {\large #2}\\[0.3em]
    \rule{\linewidth}{0.6pt}
  \end{center}
  \vspace{0.5em}
}

\newenvironment{questao}[1]{\vspace{1em}\noindent\textbf{Questão #1.}\hspace{0.4em}}{\par}
\newenvironment{solucao}{\vspace{0.3em}\noindent\textit{Solução.}\hspace{0.4em}\color{black!85}}{\par\vspace{0.5em}}


\usepackage[utf8]{inputenc}
\usepackage[T1]{fontenc}
\usepackage[brazil]{babel}
\usepackage{fullpage}
\usepackage{amsmath,amssymb}
\usepackage{enumitem}
\usepackage{xcolor}
\usepackage{fancyhdr}
\usepackage{titlesec}
\usepackage{listings}
\usepackage{hyperref}
\usepackage{booktabs}
\usepackage{graphicx}

\definecolor{matd48azul}{HTML}{0B4F6C}
\definecolor{matd48verde}{HTML}{2E8B57}

\titleformat{\section}{\normalfont\Large\bfseries\color{matd48azul}}{\thesection}{1em}{}
\titleformat{\subsection}{\normalfont\large\bfseries\color{matd48azul}}{\thesubsection}{1em}{}

\providecommand{\ListaTitulo}{Lista de Exercícios}

\setlength{\headheight}{14pt}
\pagestyle{fancy}
\fancyhf{}
\lhead{\small MATD48 --- Planejamento de Experimentos A}
\rhead{\small \ListaTitulo}
\cfoot{\thepage}
\renewcommand{\headrulewidth}{0.4pt}

\lstset{
  basicstyle=\ttfamily\small,
  breaklines=true,
  frame=single,
  columns=fullflexible,
  backgroundcolor=\color{gray!8},
  commentstyle=\color{matd48verde},
  keywordstyle=\color{matd48azul}\bfseries,
  language=R,
  extendedchars=true,
  literate=%
    {á}{{\'a}}1 {à}{{\`a}}1 {â}{{\^a}}1 {ã}{{\~a}}1
    {é}{{\'e}}1 {ê}{{\^e}}1 {í}{{\'i}}1 {ó}{{\'o}}1
    {ô}{{\^o}}1 {õ}{{\~o}}1 {ú}{{\'u}}1 {ç}{{\c c}}1
    {Á}{{\'A}}1 {À}{{\`A}}1 {Â}{{\^A}}1 {Ã}{{\~A}}1
    {É}{{\'E}}1 {Ê}{{\^E}}1 {Í}{{\'I}}1 {Ó}{{\'O}}1
    {Ô}{{\^O}}1 {Õ}{{\~O}}1 {Ú}{{\'U}}1 {Ç}{{\c C}}1
}

\hypersetup{colorlinks=true, linkcolor=matd48azul, urlcolor=matd48azul, citecolor=matd48azul}

\newcommand{\cabecalho}[2]{%
  \begin{center}
    {\Large\bfseries\color{matd48azul} #1}\\[0.3em]
    {\large #2}\\[0.3em]
    \rule{\linewidth}{0.6pt}
  \end{center}
  \vspace{0.5em}
}

\newenvironment{questao}[1]{\vspace{1em}\noindent\textbf{Questão #1.}\hspace{0.4em}}{\par}
\newenvironment{solucao}{\vspace{0.3em}\noindent\textit{Solução.}\hspace{0.4em}\color{black!85}}{\par\vspace{0.5em}}


\usepackage[utf8]{inputenc}
\usepackage[T1]{fontenc}
\usepackage[brazil]{babel}
\usepackage{fullpage}
\usepackage{amsmath,amssymb}
\usepackage{enumitem}
\usepackage{xcolor}
\usepackage{fancyhdr}
\usepackage{titlesec}
\usepackage{listings}
\usepackage{hyperref}
\usepackage{booktabs}
\usepackage{graphicx}

\definecolor{matd48azul}{HTML}{0B4F6C}
\definecolor{matd48verde}{HTML}{2E8B57}

\titleformat{\section}{\normalfont\Large\bfseries\color{matd48azul}}{\thesection}{1em}{}
\titleformat{\subsection}{\normalfont\large\bfseries\color{matd48azul}}{\thesubsection}{1em}{}

\providecommand{\ListaTitulo}{Lista de Exercícios}

\setlength{\headheight}{14pt}
\pagestyle{fancy}
\fancyhf{}
\lhead{\small MATD48 --- Planejamento de Experimentos A}
\rhead{\small \ListaTitulo}
\cfoot{\thepage}
\renewcommand{\headrulewidth}{0.4pt}

\lstset{
  basicstyle=\ttfamily\small,
  breaklines=true,
  frame=single,
  columns=fullflexible,
  backgroundcolor=\color{gray!8},
  commentstyle=\color{matd48verde},
  keywordstyle=\color{matd48azul}\bfseries,
  language=R,
  extendedchars=true,
  literate=%
    {á}{{\'a}}1 {à}{{\`a}}1 {â}{{\^a}}1 {ã}{{\~a}}1
    {é}{{\'e}}1 {ê}{{\^e}}1 {í}{{\'i}}1 {ó}{{\'o}}1
    {ô}{{\^o}}1 {õ}{{\~o}}1 {ú}{{\'u}}1 {ç}{{\c c}}1
    {Á}{{\'A}}1 {À}{{\`A}}1 {Â}{{\^A}}1 {Ã}{{\~A}}1
    {É}{{\'E}}1 {Ê}{{\^E}}1 {Í}{{\'I}}1 {Ó}{{\'O}}1
    {Ô}{{\^O}}1 {Õ}{{\~O}}1 {Ú}{{\'U}}1 {Ç}{{\c C}}1
}

\hypersetup{colorlinks=true, linkcolor=matd48azul, urlcolor=matd48azul, citecolor=matd48azul}

\newcommand{\cabecalho}[2]{%
  \begin{center}
    {\Large\bfseries\color{matd48azul} #1}\\[0.3em]
    {\large #2}\\[0.3em]
    \rule{\linewidth}{0.6pt}
  \end{center}
  \vspace{0.5em}
}

\newenvironment{questao}[1]{\vspace{1em}\noindent\textbf{Questão #1.}\hspace{0.4em}}{\par}
\newenvironment{solucao}{\vspace{0.3em}\noindent\textit{Solução.}\hspace{0.4em}\color{black!85}}{\par\vspace{0.5em}}


\begin{document}

\cabecalho{Gabarito --- Lista de Exercícios 02}{Modelo linear, mínimos quadrados e estimabilidade (Capítulo 2 / Aula 02)}

\begin{questao}{1} \textbf{(Ciência de dados --- mínimos quadrados na mão)}
\end{questao}
\begin{solucao}
\begin{enumerate}[label=(\alph*)]
  \item
  $$
  \mathbf{X} = \begin{bmatrix} 1 & 2 \\ 1 & 4 \\ 1 & 6 \\ 1 & 8 \\ 1 & 10 \end{bmatrix}, \qquad
  \mathbf{Y} = \begin{bmatrix} 15 \\ 22 \\ 25 \\ 33 \\ 38 \end{bmatrix}.
  $$
  \item
  $$
  \mathbf{X}'\mathbf{X} = \begin{bmatrix} 5 & 30 \\ 30 & 220 \end{bmatrix}, \qquad
  \mathbf{X}'\mathbf{Y} = \begin{bmatrix} 133 \\ 912 \end{bmatrix}.
  $$
  \item Com $a=5$, $b=30$, $c=220$: $ac - b^2 = 1100 - 900 = 200$, então
  $$
  (\mathbf{X}'\mathbf{X})^{-1} = \frac{1}{200}\begin{bmatrix} 220 & -30 \\ -30 & 5 \end{bmatrix},
  $$
  $$
  \hat{\boldsymbol\beta} = (\mathbf{X}'\mathbf{X})^{-1}\mathbf{X}'\mathbf{Y} =
  \frac{1}{200}\begin{bmatrix} 220(133) - 30(912) \\ -30(133) + 5(912) \end{bmatrix} =
  \frac{1}{200}\begin{bmatrix} 1900 \\ 570 \end{bmatrix} = \begin{bmatrix} 9.5 \\ 2.85 \end{bmatrix}.
  $$
  Logo $\hat\beta_0 = 9.5$ e $\hat\beta_1 = 2.85$.
  \item Cada script adicional carregado está associado a um aumento estimado de \textbf{2,85
    décimos de segundo} (0,285 s) no tempo de carregamento da página, mantendo tudo o mais igual.
  \item A soma dos resíduos deve ser \textbf{zero}. Das equações normais,
    $\mathbf{X}'\mathbf{e} = \mathbf{X}'(\mathbf{Y}-\mathbf{X}\hat{\boldsymbol\beta}) = \mathbf{0}$.
    A primeira linha de $\mathbf{X}'$ é um vetor de 1's (porque a primeira coluna de $\mathbf{X}$
    é o intercepto), então a primeira componente de $\mathbf{X}'\mathbf{e}$ é exatamente
    $\sum_i e_i$, que portanto é zero --- sem precisar recalcular os resíduos individualmente.
    (Conferindo: $\hat{\mathbf{Y}} = (15.2,\ 20.9,\ 26.6,\ 32.3,\ 38.0)$, $\mathbf{e} = (-0.2,\
    1.1,\ -1.6,\ 0.7,\ 0.0)$, soma $= 0$.)
\end{enumerate}
\end{solucao}

\begin{questao}{2} \textbf{(Psicologia --- estimabilidade)}
\end{questao}
\begin{solucao}
\begin{enumerate}[label=(\alph*)]
  \item
  $$
  \mathbf{X} = \begin{bmatrix}
  1 & 1 & 0 & 0 \\ 1 & 1 & 0 & 0 \\
  1 & 0 & 1 & 0 \\ 1 & 0 & 1 & 0 \\
  1 & 0 & 0 & 1 \\ 1 & 0 & 0 & 1
  \end{bmatrix}.
  $$
  \item O posto é \textbf{3}, não 4. A coluna de $\mu$ é exatamente a soma das colunas de
    $\alpha_1$, $\alpha_2$ e $\alpha_3$ (cada linha tem um único $\alpha_j$ igual a 1, então a
    soma das três colunas de $\alpha$ é uma coluna de 1's, idêntica à coluna de $\mu$) --- uma
    dependência linear exata que reduz o posto em 1.
  \item Usando o critério $\boldsymbol\lambda'\boldsymbol\beta$ estimável $\iff \boldsymbol\lambda'$
    no espaço-linha de $\mathbf{X}$ (equivalentemente, existe $\mathbf{a}$ com
    $\mathbf{a}'\mathbf{X}=\boldsymbol\lambda'$):
    \begin{enumerate}[label=(\roman*)]
      \item $\mu$: \textbf{não estimável}. Não há combinação das seis linhas de $\mathbf{X}$ que
        produza $(1,0,0,0)$ --- toda linha de $\mathbf{X}$ tem um $\mu$ e um $\alpha_j$
        simultaneamente "ligados"; isolar $\mu$ exigiria cancelar todos os $\alpha_j$'s ao mesmo
        tempo em que se mantém o coeficiente de $\mu$ em 1, o que não é possível dado o
        posto 3 da matriz.
      \item $\mu + \alpha_1$: \textbf{estimável}. Tome $\mathbf{a} = (\tfrac12, \tfrac12, 0,0,0,0)$
        (a média das duas observações do grupo 1): $\mathbf{a}'\mathbf{X} = (1,1,0,0)$, que é
        exatamente $\boldsymbol\lambda'$ para $\mu+\alpha_1$. É a média populacional do grupo 1.
      \item $\alpha_1 - \alpha_2$: \textbf{estimável}. Tome $\mathbf{a} = (\tfrac12,\tfrac12,
        -\tfrac12,-\tfrac12,0,0)$: $\mathbf{a}'\mathbf{X} = (0,1,-1,0)$, a diferença entre as
        médias dos grupos 1 e 2 --- um contraste, sempre estimável independentemente da
        parametrização.
      \item $\alpha_1+\alpha_2+\alpha_3$: \textbf{não estimável}. Não existe $\mathbf{a}$ com
        $\mathbf{a}'\mathbf{X} = (0,1,1,1)$: qualquer combinação das linhas que dê coeficiente 1
        em cada $\alpha_j$ necessariamente também carrega um coeficiente não nulo em $\mu$ (soma
        dos pesos $\mathbf{a}$), então essa função particular não pode ser isolada.
    \end{enumerate}
\end{enumerate}
\end{solucao}

\begin{questao}{3} \textbf{(Agricultura --- especificação do modelo)}
\end{questao}
\begin{solucao}
\begin{enumerate}[label=(\alph*)]
  \item Ordenando as parcelas (baixa,V1), (baixa,V2), (alta,V1), (alta,V2):
  $$
  \mathbf{X} = \begin{bmatrix} 1 & 0 & 0 \\ 1 & 0 & 1 \\ 1 & 1 & 0 \\ 1 & 1 & 1 \end{bmatrix}.
  $$
  \item Sim, $\mathbf{X}$ tem posto coluna completo (posto 3 = número de colunas). As quatro
    linhas são distintas e, tomando três delas --- por exemplo, as linhas 1, 2 e 3 --- obtém-se
    uma submatriz $3\times 3$ triangular inferior com diagonal $(1,1,1)$ e determinante não nulo,
    o que já basta para garantir posto 3 completo em $\mathbf{X}$.
  \item $\beta_0$ é a produtividade média esperada da combinação de referência (dose baixa,
    V1). $\beta_1$ é a diferença esperada na produtividade média entre dose alta e dose baixa,
    mantendo a variedade fixa (em V1). $\beta_2$ é a diferença esperada entre V2 e V1, mantendo a
    dose fixa (em baixa). Como não há termo de interação neste modelo, assume-se que essas duas
    diferenças são as mesmas independentemente do nível do outro fator.
\end{enumerate}
\end{solucao}

\begin{questao}{4} \textbf{(Ciência de dados --- geometria dos mínimos quadrados)}
\end{questao}
\begin{solucao}
\begin{enumerate}[label=(\alph*)]
  \item Das equações normais, $\mathbf{X}'\mathbf{X}\hat{\boldsymbol\beta} = \mathbf{X}'\mathbf{Y}$,
    logo $\mathbf{X}'\mathbf{X}\hat{\boldsymbol\beta} - \mathbf{X}'\mathbf{Y} = \mathbf{0}$, ou
    seja, $-\mathbf{X}'(\mathbf{Y} - \mathbf{X}\hat{\boldsymbol\beta}) = \mathbf{0}$. Como
    $\mathbf{e} = \mathbf{Y}-\mathbf{X}\hat{\boldsymbol\beta}$, isso é exatamente
    $\mathbf{X}'\mathbf{e} = \mathbf{0}$.
  \item Se a primeira coluna de $\mathbf{X}$ é $\mathbf{1}$ (vetor de 1's), a primeira linha de
    $\mathbf{X}'$ também é $\mathbf{1}'$, e a primeira componente de $\mathbf{X}'\mathbf{e}$ é
    $\mathbf{1}'\mathbf{e} = \sum_i e_i$. Pelo item (a), essa componente é zero, logo $\sum_i e_i
    = 0$.
  \item Não. $\sum_i e_i = 0$ é uma restrição \emph{global} (média dos resíduos zero); ela não
    impede padrões \emph{locais} sistemáticos. Por exemplo, os resíduos poderiam ser
    sistematicamente positivos para valores pequenos de $x$ e sistematicamente negativos para
    valores grandes de $x$ (indicando uma relação não linear que um modelo linear não capturou) e
    ainda assim somar exatamente zero. A propriedade garante apenas ausência de viés
    \emph{médio}, não ausência de má especificação do modelo --- por isso sempre inspecionamos o
    gráfico de resíduos, não só sua soma.
\end{enumerate}
\end{solucao}

\begin{questao}{5} \textbf{(Ciência de dados --- exercício computacional em R)}
\end{questao}
\begin{solucao}
\begin{enumerate}[label=(\alph*)]
  \item
\begin{lstlisting}
x <- c(2, 4, 6, 8, 10)
y <- c(15, 22, 25, 33, 38)
X <- cbind(1, x)

beta_hat <- solve(t(X) %*% X) %*% t(X) %*% y

y_hat <- X %*% beta_hat
residuos <- y - y_hat

sum(residuos)
\end{lstlisting}
  \item Rodando o código, \texttt{beta\_hat} retorna $(9.5,\ 2.85)$ --- coincide exatamente com o
    resultado manual da Questão 1(c), e \texttt{sum(residuos)} retorna um valor numericamente
    igual a zero (algo como \texttt{-2.1e-13}, um artefato de precisão de ponto flutuante).
  \item
\begin{lstlisting}
coef(lm(y ~ x))
\end{lstlisting}
    As duas estimativas coincidem porque \texttt{lm()} resolve exatamente as mesmas equações
    normais $\mathbf{X}'\mathbf{X}\hat{\boldsymbol\beta}=\mathbf{X}'\mathbf{Y}$ (internamente via
    decomposição QR, por estabilidade numérica, mas o resultado algébrico é idêntico ao da
    inversão direta de $\mathbf{X}'\mathbf{X}$).
\end{enumerate}
\end{solucao}

\begin{questao}{6} \textbf{(Dedução --- equações normais)}
\end{questao}
\begin{solucao}
\begin{enumerate}[label=(\alph*)]
  \item
  $$
  S(\boldsymbol\beta) = (\mathbf{Y}-\mathbf{X}\boldsymbol\beta)'(\mathbf{Y}-\mathbf{X}\boldsymbol\beta)
  = \mathbf{Y}'\mathbf{Y} - \mathbf{Y}'\mathbf{X}\boldsymbol\beta - \boldsymbol\beta'\mathbf{X}'\mathbf{Y}
  + \boldsymbol\beta'\mathbf{X}'\mathbf{X}\boldsymbol\beta.
  $$
  Como $\mathbf{Y}'\mathbf{X}\boldsymbol\beta$ é um escalar, ele é igual à sua própria transposta
  $\boldsymbol\beta'\mathbf{X}'\mathbf{Y}$, então os dois termos do meio se combinam:
  $$
  S(\boldsymbol\beta) = \mathbf{Y}'\mathbf{Y} - 2\boldsymbol\beta'\mathbf{X}'\mathbf{Y} +
  \boldsymbol\beta'\mathbf{X}'\mathbf{X}\boldsymbol\beta.
  $$
  \item Aplicando as duas regras de derivação termo a termo, com $\mathbf{a} =
  \mathbf{X}'\mathbf{Y}$ no termo linear e $\mathbf{A} = \mathbf{X}'\mathbf{X}$ (simétrica) no
  termo quadrático:
  $$
  \frac{\partial S(\boldsymbol\beta)}{\partial \boldsymbol\beta} = -2\mathbf{X}'\mathbf{Y} +
  2\mathbf{X}'\mathbf{X}\boldsymbol\beta.
  $$
  \item Igualando a $\mathbf{0}$:
  $$
  -2\mathbf{X}'\mathbf{Y} + 2\mathbf{X}'\mathbf{X}\hat{\boldsymbol\beta} = \mathbf{0}
  \ \Longrightarrow\ \mathbf{X}'\mathbf{X}\hat{\boldsymbol\beta} = \mathbf{X}'\mathbf{Y}.
  $$
  \item A Hessiana de $S(\boldsymbol\beta)$ é a matriz de segundas derivadas,
  $\partial^2 S/\partial\boldsymbol\beta\partial\boldsymbol\beta' = 2\mathbf{X}'\mathbf{X}$. Para
  qualquer vetor $\mathbf{v}$, $\mathbf{v}'(2\mathbf{X}'\mathbf{X})\mathbf{v} =
  2(\mathbf{X}\mathbf{v})'(\mathbf{X}\mathbf{v}) = 2\lVert \mathbf{X}\mathbf{v}\rVert^2 \geq 0$,
  logo a Hessiana é semidefinida positiva em qualquer ponto do domínio (não depende de
  $\boldsymbol\beta$). Uma função cuja Hessiana é semidefinida positiva em todo o domínio é
  \textbf{convexa}; para uma função convexa, todo ponto crítico (derivada primeira nula) é um
  \textbf{mínimo global}, nunca um máximo ou ponto de sela. Isso garante que qualquer solução das
  equações normais realmente minimiza $S(\boldsymbol\beta)$, e não apenas anula sua derivada.
\end{enumerate}
\end{solucao}

\begin{questao}{7} \textbf{(Dedução/computacional --- pseudoinversa de Moore-Penrose via SVD)}
\end{questao}
\begin{solucao}
\begin{enumerate}[label=(\alph*)]
  \item
  $$
  \mathbf{A}\mathbf{v}_3 = \begin{bmatrix}4&2&2\\2&2&0\\2&0&2\end{bmatrix}
  \frac{1}{\sqrt3}\begin{bmatrix}1\\-1\\-1\end{bmatrix} =
  \frac{1}{\sqrt3}\begin{bmatrix}4-2-2\\2-2-0\\2-0-2\end{bmatrix} =
  \frac{1}{\sqrt3}\begin{bmatrix}0\\0\\0\end{bmatrix} = \mathbf{0}.
  $$
  Confirmado: $\mathbf{v}_3$ é autovetor de autovalor $0$. Isso já era esperado porque
  $\mathbf{X}$ (a matriz de delineamento original) tem posto $2$ (a coluna de $\mu$ é a soma das
  colunas de $\tau_1,\tau_2$), então $\mathbf{A}=\mathbf{X}'\mathbf{X}$ também tem posto $2$ (o
  posto de $\mathbf{X}'\mathbf{X}$ é sempre igual ao posto de $\mathbf{X}$) --- uma matriz
  $3\times3$ de posto $2$ é singular e tem exatamente um autovalor nulo.
  \item
  $$
  \mathbf{D}^+ = \begin{bmatrix} 1/6 & 0 & 0 \\ 0 & 1/2 & 0 \\ 0 & 0 & 0 \end{bmatrix}.
  $$
  \item Como o termo em $\mathbf{v}_3$ é multiplicado por $0$ em $\mathbf{D}^+$,
  $$
  \mathbf{A}^+ = \tfrac16 \mathbf{v}_1\mathbf{v}_1' + \tfrac12 \mathbf{v}_2\mathbf{v}_2'.
  $$
  A entrada $(1,1)$ usa só as primeiras componentes: $v_{1,1} = 2/\sqrt6$, $v_{2,1}=0$. Logo
  $$
  [\mathbf{A}^+]_{11} = \tfrac16 (2/\sqrt6)^2 + \tfrac12 (0)^2 = \tfrac16 \cdot \tfrac46 =
  \tfrac{4}{36} = \tfrac19 \approx 0{,}1111,
  $$
  que bate exatamente com a entrada $(1,1)$ de \texttt{MASS::ginv(A)}.
  \item
\begin{lstlisting}
A <- matrix(c(4,2,2, 2,2,0, 2,0,2), ncol = 3, byrow = TRUE)
sv <- svd(A)
D_mais <- diag(ifelse(sv$d > 1e-8, 1 / sv$d, 0))
A_mais <- sv$v %*% D_mais %*% t(sv$u)
max(abs(A_mais - MASS::ginv(A)))
\end{lstlisting}
  Rodando o código, a diferença máxima é numericamente zero: a construção manual via SVD e
  \texttt{MASS::ginv()} produzem a mesma matriz.
  \item
  $$
  \mathbf{b} = \mathbf{A}^+\mathbf{X}'\mathbf{y} \approx (9{,}5833,\ 5{,}1667,\ 4{,}4167)',
  $$
  e as funções estimáveis:
  $$
  \mu+\tau_1 = 9{,}5833+5{,}1667 = 14{,}75, \qquad \mu+\tau_2 = 9{,}5833+4{,}4167 = 14{,}00,
  \qquad \tau_1-\tau_2 = 5{,}1667-4{,}4167 = 0{,}75.
  $$
  Esses três valores coincidem com os que qualquer outra inversa generalizada de $\mathbf{A}$
  produziria (mesmo que o vetor $\mathbf{b}$ intermediário seja diferente), porque
  $\mu+\tau_1$, $\mu+\tau_2$ e $\tau_1-\tau_2$ são funções \textbf{estimáveis} --- e o texto do
  Capítulo 2 do livro mostra que o valor de uma função estimável $\boldsymbol\lambda'\mathbf{b}$
  não depende de qual inversa generalizada $(\mathbf{X}'\mathbf{X})^-$ foi usada para obter
  $\mathbf{b}$, só o próprio $\mathbf{b}$ depende dessa escolha.
\end{enumerate}
\end{solucao}

\end{document}
